% ============================================================
%  Growing, Coasting, or Stuck: Engineer Career Trajectories
%  and Communication Patterns in Open-Source Data
%  Draft for internal review — not for distribution
% ============================================================
\documentclass[12pt,a4paper]{article}

\usepackage[margin=2.5cm]{geometry}
\usepackage{amsmath,amssymb}
\usepackage{booktabs}
\usepackage{hyperref}
\usepackage{graphicx}
\usepackage{array}
\usepackage{parskip}
\usepackage{setspace}
\usepackage{caption}
\usepackage[numbers]{natbib}
\usepackage{xcolor}
\usepackage{titlesec}

\onehalfspacing
\setlength{\parindent}{0pt}
\setlength{\parskip}{8pt}

\title{%
  \textbf{Growing, Coasting, or Stuck:}\\
  Engineer Career Trajectories and Communication Patterns\\
  in Open-Source Software Projects\\[6pt]
  \large\textit{Working paper — internal review draft}
}
\author{%
  Task2Vec Research\\
  \small\texttt{research@task2vec.com}
}
\date{March 2026}

\begin{document}
\maketitle
\thispagestyle{empty}

% ── Abstract ────────────────────────────────────────────────────────
\begin{abstract}
The 70-20-10 learning framework posits that 70\,\% of professional
development occurs through on-the-job experience.
We test this assumption empirically using 45,631 Jira tickets
linked to git commit histories across three major Apache Software
Foundation projects: Apache Camel, Apache Spark, and Apache Hadoop,
covering 378 individual contributors over a 15-year window (2007–2022).

We define a composite \emph{complexity score} for each resolved ticket
based on the number of software modules touched, file count (log-scaled),
and involvement of architecturally central (``core'') modules.
For each contributor with at least 10 linked tickets, we compute a
Spearman rank correlation between ticket index and complexity score,
yielding a career arc signal.

The distribution of career arcs directly contradicts the 70-20-10
assumption: only \textbf{11\,\%} of contributors exhibit a
\emph{growing} trajectory (rising complexity, improving velocity),
\textbf{20\,\%} exhibit a \emph{coasting} trajectory (declining
complexity, slowing velocity), and \textbf{70\,\%} show no measurable
directional effect.

A secondary analysis of the Apache Spark communication network
(co-comment graph, $n=372$ contributors with $\geq20$ project
comments) reveals a striking inversion: contributors on a coasting
trajectory have, on average, 13$\times$ more comments and 3$\times$
higher degree centrality than growing contributors.
This suggests that visibility and communication volume in a
community is not a reliable proxy for career development.

A multi-factor predictor analysis (Mann-Whitney $U$, 21 variables)
identifies sub-task ratio and scope delta as the strongest actionable
predictors of a growing trajectory, and shows that Apache Hadoop's
uniformly hard ticket environment --- despite having 6--18$\times$
higher median complexity than Camel or Spark --- produces
\textit{no statistically distinguishable} career signal.
We term this the \textit{complexity floor effect}: environmental
difficulty is necessary but not sufficient for growth.

We discuss implications for engineering talent management,
the design of deliberate work-assignment curricula, and the
limits of generalising from open-source to commercial settings.
All analysis is reproducible from publicly available data.
\end{abstract}

\newpage
\tableofcontents
\newpage

% ── 1. Introduction ─────────────────────────────────────────────────
\section{Introduction}

The 70-20-10 model, first articulated by McCall, Lombardo, and
Eichinger \citep{McCall1988}, proposes that professional learning
derives 70\,\% from on-the-job experience, 20\,\% from social
learning (feedback, mentorship, peer interaction), and 10\,\% from
formal instruction. Despite limited empirical validation
\citep{Clardy2018}, the framework has become a dominant organising
principle in corporate learning and development programmes.

For software engineers specifically, the model implies that sustained
participation in a software project --- resolving tickets, writing
code, reviewing pull requests --- should, over time, expand the
range and complexity of work an engineer can handle. The implicit
mechanism is deliberate practice through exposure: repeated problem
solving, cognitive load, and peer feedback through code review and
commenting.

We ask a simple empirical question: \textit{does the available
evidence from large-scale open-source project histories support this
assumption?}

Open-source software projects offer an unusual opportunity for
this inquiry. The Apache Software Foundation maintains public Jira
archives and git repositories in which individual contributor
identities, ticket metadata, and code-change records are all
accessible and linkable. Unlike commercial settings, where HR data
is proprietary and ticket quality varies by team, Apache projects
apply consistent issue-tracking conventions across decades.

Our contribution is threefold:
\begin{enumerate}
  \item We introduce a \textit{career arc} measure derived entirely
    from Jira and git metadata --- requiring no surveys, performance
    reviews, or instrumentation beyond standard engineering tooling.
  \item We find that the 70-20-10 assumption does not hold empirically:
    the majority of contributors show no measurable career-arc signal,
    and the coasting group is nearly twice the size of the growing group.
  \item We find that coasting contributors dominate communication
    networks by volume and centrality, suggesting that managerial
    ``visibility'' is inversely correlated with growth trajectory.
  \item We identify the two most actionable predictors of a growing
    trajectory: a low sub-task assignment ratio and positive scope
    delta (rising files-per-commit over the career).
    We additionally show that a uniformly hard ticket environment
    --- the intuitive intervention for accelerating growth --- can
    paradoxically mask the career signal by compressing within-engineer
    variance.
\end{enumerate}

The remainder of the paper is structured as follows. Section~2
describes the dataset and preprocessing steps. Section~3 defines
the complexity score and career arc methodology. Section~4 presents
the communication network analysis. Section~5 reports results.
Section~6 discusses implications and limitations. Section~7 concludes.

% ── 2. Data ─────────────────────────────────────────────────────────
\section{Data}

\subsection{Sources}

We use three Apache Software Foundation projects:

\begin{itemize}
  \item \textbf{Apache Camel} (integration framework): 14,184 tickets
    matched to commits; 85\,\% coverage of the full Jira archive.
  \item \textbf{Apache Spark} (distributed data processing): 24,994
    tickets matched; 71\,\% coverage.
  \item \textbf{Apache Hadoop} (distributed storage): 6,453 tickets
    matched; 49\,\% coverage.
\end{itemize}

The lower coverage for Hadoop is attributable to inconsistent
commit message conventions in the early project history (pre-2012).

Jira tickets were retrieved via the public REST API endpoint
\texttt{https://issues.apache.org/jira/\\rest/api/2/search}.
Git repositories were cloned with \texttt{--filter=blob:none} to
minimise download size. Only \textit{resolved} tickets --- those
with a non-null \texttt{resolutiondate} field --- were included.

\subsection{Linking tickets to commits}

Each repository commit message was scanned for references matching
the pattern \texttt{PROJECT-[0-9]+} (e.g.\ \texttt{SPARK-12345}).
The changed files for each matched commit were extracted using
\texttt{git diff-tree --no-commit-id -r}.
Where multiple commits referenced the same ticket key, all changed
files were pooled.

\subsection{Contributor identification}

Apache Jira anonymises contributor identities as UUIDs
(e.g.\ \texttt{author\_key|4e4cf22c-...|}).
We treat each unique UUID as a distinct contributor.
We cannot confirm identity stability across accounts,
though Jira UUIDs are stable within the system lifetime.

\subsection{Dataset summary}

Table~\ref{tab:dataset} summarises the final dataset used in the
career trajectory analysis.

\begin{table}[h]
\centering
\caption{Dataset summary}
\label{tab:dataset}
\begin{tabular}{lrr}
\toprule
Statistic & Value \\
\midrule
Total linked tickets (all projects) & 45,631 \\
Unique contributors (all projects) & 11,185 \\
Contributors with $\geq10$ linked tickets & 378 \\
History window & 2007–2022 (15 years) \\
Median tickets per qualifying contributor & 54 \\
Median career span (qualifying contributors) & 4.1 years \\
\bottomrule
\end{tabular}
\end{table}

% ── 3. Methods: Complexity and Career Arc ───────────────────────────
\section{Complexity Score and Career Arc Methodology}

\subsection{Ticket complexity score}

For each resolved ticket $t$, let:
\begin{itemize}
  \item $M_t$ = number of distinct software modules touched by the
    associated commits (e.g.\ \texttt{camel-core}, \texttt{spark-sql}).
  \item $F_t$ = total number of files changed across all associated commits.
  \item $C_t \in [0,1]$ = fraction of changed files belonging to
    \textit{core modules} (defined per project; see below).
\end{itemize}

The raw complexity score is:
\begin{equation}
  \text{complexity\_raw}(t) \;=\;
  0.4 \cdot M_t \;+\; 0.3 \cdot \log(F_t + 1) \;+\; 0.3 \cdot 5\,C_t
  \label{eq:complexity_raw}
\end{equation}

Scores are then min-max normalised across all tickets in the dataset:
\begin{equation}
  \text{complexity}(t) \;=\;
  \frac{\text{complexity\_raw}(t) - \min_s(\text{complexity\_raw}(s))}
       {\max_s(\text{complexity\_raw}(s)) - \min_s(\text{complexity\_raw}(s))}
  \label{eq:complexity_norm}
\end{equation}

\textbf{Core module identification.} Core modules were identified by
examining which modules appear most frequently among tickets in the
top quartile of watcher count and the slowest quartile of resolution
time --- a proxy for architectural centrality and community-perceived
difficulty. For Apache Camel, core modules are \texttt{camel-core},
\texttt{camel-api}, and \texttt{camel-support}. For Spark:
\texttt{core} and \texttt{common}. For Hadoop:
\texttt{hadoop-common} and \texttt{hadoop-auth}.

\textbf{Rationale for weights.} Module count (weight 0.4) is the
dominant term because touching multiple modules requires
cross-system understanding. File count is log-scaled and
weighted 0.3 to capture scope while limiting the influence of
large but mechanically simple refactors. Core module fraction
(weight 0.3) captures architectural centrality.
Weights were selected by qualitative review of boundary cases
and held fixed; they were not tuned to optimise any
downstream measure.

\textbf{Validation.} High-complexity tickets (top quartile by score)
resolved in a median of 8 days vs.\ 2 days for low-complexity tickets,
and attracted a median of 5 comments vs.\ 2 --- suggesting the score
captures genuine difficulty and not merely code volume.

\subsection{Career arc: Spearman rank correlation}

For each qualifying contributor $u$ (with $\geq10$ linked tickets),
let $(t_1, t_2, \ldots, t_n)$ be their tickets ordered by
\texttt{created} date, and let $c_i = \text{complexity}(t_i)$.

The \textit{complexity arc} is the Spearman rank correlation
between ticket index and complexity:
\begin{equation}
  \rho_{\text{complexity}}(u) \;=\;
  r_s\bigl(\,(1,2,\ldots,n),\,(c_1, c_2, \ldots, c_n)\,\bigr)
  \label{eq:rho_complexity}
\end{equation}
where $r_s$ denotes the Spearman rank correlation coefficient.

The \textit{velocity arc} is analogously defined as the Spearman
correlation between ticket index and resolution time in days
$d_i = \text{days}(t_i)$:
\begin{equation}
  \rho_{\text{velocity}}(u) \;=\;
  r_s\bigl(\,(1,2,\ldots,n),\,(d_1, d_2, \ldots, d_n)\,\bigr)
  \label{eq:rho_velocity}
\end{equation}

\textbf{Trajectory labels.} Using thresholds of $|\rho| = 0.2$,
contributors are assigned to one of three groups:
\begin{align}
  \text{Growing}   &: \quad \rho_\text{complexity} > +0.2
    \;\text{ AND }\; \rho_\text{velocity} < -0.2 \label{eq:growing} \\
  \text{Coasting}  &: \quad \rho_\text{complexity} < -0.2
    \;\text{ AND }\; \rho_\text{velocity} > +0.2 \label{eq:coasting} \\
  \text{No change} &: \quad \text{otherwise} \label{eq:nochange}
\end{align}

The threshold of 0.2 is the conventional ``weak correlation''
boundary in the Spearman literature.
We report robustness checks at thresholds of 0.15 and 0.25 in
Section~\ref{sec:robustness}.

\textbf{Why Spearman?} We use rank correlation rather than
Pearson to avoid sensitivity to outlier tickets (a single
very complex ticket in year 1 or year 10 could dominate a
linear regression). Spearman measures the monotonic trend,
which is the appropriate estimand for ``is this contributor
consistently taking on harder or easier work over time?''

% ── 4. Communication Network ────────────────────────────────────────
\section{Communication Network Analysis}

\subsection{Co-comment graph construction}

For each Jira ticket $k$ in the Apache Spark project, let
$A(k) \subseteq U$ denote the set of contributors who left at
least one comment on that ticket.
The \textit{co-comment graph} $G = (V, E, w)$ is defined as:
\begin{align}
  V &= \{u \in U : \text{comment\_count}(u) \geq 20\} \\
  E &= \bigl\{\{u,v\} : \exists\, k \text{ s.t. } u,v \in A(k)\bigr\} \\
  w(u,v) &= \bigl|\{k : u \in A(k) \text{ and } v \in A(k)\}\bigr|
\end{align}

Only edges with $w(u,v) \geq 2$ are retained.
The largest connected component (LCC) is used in all analyses.

\subsection{Network metrics}

For the LCC with $n$ nodes and $m$ edges:

\textbf{Density:}
\begin{equation}
  \delta(G) = \frac{2m}{n(n-1)}
\end{equation}

\textbf{Degree centrality} of node $u$:
\begin{equation}
  C_D(u) = \frac{\deg(u)}{n-1}
\end{equation}

\textbf{Betweenness centrality} of node $u$:
\begin{equation}
  C_B(u) = \frac{1}{(n-1)(n-2)}
  \sum_{s \neq t \neq u}
  \frac{\sigma_{st}(u)}{\sigma_{st}}
\end{equation}
where $\sigma_{st}$ is the number of shortest paths from $s$ to $t$
and $\sigma_{st}(u)$ is the number passing through $u$.
Weighted shortest paths use edge weight $w(u,v)$.

\textbf{Community detection} was performed using the greedy
modularity maximisation algorithm \citep{Clauset2004}, which
iteratively merges communities to maximise:
\begin{equation}
  Q = \frac{1}{2m} \sum_{u,v}
  \left[w(u,v) - \frac{k_u k_v}{2m}\right]
  \delta(c_u, c_v)
\end{equation}
where $k_u = \sum_v w(u,v)$ is the weighted degree of node $u$,
$c_u$ is the community assignment of $u$, and $\delta$ is the
Kronecker delta.

\subsection{Visualisation}

The co-comment graph was visualised using a force-directed
(Fruchterman-Reingold) spring layout \citep{Fruchterman1991}
with edge weights inverted so that more frequently co-commenting
pairs are positioned closer:
\begin{equation}
  \text{spring\_length}(u,v) \propto \frac{1}{w(u,v)^{0.6}}
\end{equation}
Node size is proportional to $\log(\text{comment\_count})$.
Node colour encodes career trajectory label.

% ── 5. Results ──────────────────────────────────────────────────────
\section{Results}

\subsection{Career trajectory distribution}

Table~\ref{tab:trajectories} shows the distribution of trajectory
labels across all 378 qualifying contributors.

\begin{table}[h]
\centering
\caption{Career trajectory distribution ($n=378$ contributors,
three Apache projects, 2007–2022)}
\label{tab:trajectories}
\begin{tabular}{lrrl}
\toprule
Trajectory & $n$ & \% & Definition \\
\midrule
Growing   & 42  & 11.1\,\% &
  $\rho_\text{complexity} > 0.2$ and $\rho_\text{velocity} < -0.2$ \\
Coasting  & 76  & 20.1\,\% &
  $\rho_\text{complexity} < -0.2$ and $\rho_\text{velocity} > 0.2$ \\
No change & 260 & 68.8\,\% & Neither condition met \\
\bottomrule
\end{tabular}
\end{table}

The distribution is robust to project: Growing ranges from 9\,\%
(Hadoop) to 13\,\% (Camel); Coasting ranges from 17\,\% (Spark)
to 24\,\% (Camel).

\textbf{The escalation arc.}
Growing contributors follow a characteristic three-phase pattern:
(1) initial repetitions in a single module cluster with declining
velocity; (2) scope expansion with simultaneously improving velocity
(the \textit{transfer signal}); (3) rising core-module involvement
and watcher counts.
Nine contributors (2.4\,\% of the total) achieved full mastery
--- rising complexity with improving velocity throughout the
entire career arc.

\subsection{The coasting paradox: communication volume vs.\ trajectory}

A key finding emerges from combining trajectory labels with
communication network data.
Table~\ref{tab:comm_stats} reports network statistics broken down
by trajectory group for the Apache Spark LCC ($n=372$).

\begin{table}[h]
\centering
\caption{Communication network statistics by trajectory group
(Apache Spark, LCC, $n=372$)}
\label{tab:comm_stats}
\begin{tabular}{lrrrr}
\toprule
Group & $n$ & Avg.\ comments & Avg.\ $C_D$ & Avg.\ $C_B$ \\
\midrule
Coasting  & 12  & 3,637 & 0.216 & 0.022 \\
No change & 77  & 573   & 0.113 & 0.013 \\
Growing   & 17  & 283   & 0.066 & 0.006 \\
No data   & 266 & 48    & 0.021 & 0.001 \\
\bottomrule
\end{tabular}
\end{table}

Coasting contributors produce, on average, \textbf{13 times more
comments} than growing contributors, and have \textbf{3.3 times
higher degree centrality}.
Betweenness centrality --- which measures how often a node lies on
the shortest path between two others --- is 3.6 times higher for
coasting contributors.

This pattern, which we term the \textit{coasting paradox}, implies
that the contributors most visible in the communication network
(highest comment volume, most connections, greatest bridging
influence) are precisely those whose career arc shows declining
complexity and slowing velocity.

Importantly, spring-layout distance from the network centroid is
nearly identical for Growing (0.282), Coasting (0.276), and No
change (0.273) groups. The visual dominance of coasting nodes in
network visualisations therefore arises from node size (comment
volume), not positional centrality. The structural difference is
in degree and betweenness, not in physical distance to core.

\subsection{Communication network structure}

The Apache Spark co-comment graph has the following properties
(LCC, $n=372$):
\begin{itemize}
  \item Nodes: 372 \quad Edges: 2{,}635 \quad Density: 0.038
  \item Greedy modularity: $Q = 0.47$, indicating strong community
    structure
  \item 9 communities detected; two dominant communities
    (C0: 70 members, C1: 61 members)
  \item Top inter-community edge weight: 17,610
    (between communities C0 and C1)
\end{itemize}

Modularity $Q = 0.47$ is well above the $Q > 0.3$ threshold
conventionally associated with meaningful community structure
\citep{Newman2004}.

\subsection{Co-change network}
\label{sec:cochange}

A secondary analysis examined structural coupling in the codebase
via file co-change pairs.
From 23,539 multi-file commits across the three projects,
325,908 file-pair co-occurrence records were extracted.

Files in the top 10\,\% of co-change frequency resolved associated
tickets 2.8$\times$ slower and attracted 3.1$\times$ more comments
than low-frequency files, providing independent validation that
co-change degree captures genuine architectural centrality.

Contributors who worked regularly on high-degree co-change files
(``hub files'') showed a higher Growing fraction (18\,\% vs.\
11\,\% overall) and a lower Coasting fraction (14\,\% vs.\ 20\,\%),
suggesting that structural exposure to architectural hubs is a
driver of career growth.

\subsection{Multi-factor growth predictors}
\label{sec:multifactor}

To understand \textit{what distinguishes growing contributors from others}
--- and why Hadoop shows a different trajectory distribution ---
we conducted a multi-factor analysis using Mann-Whitney $U$ tests
across 21 candidate predictors.
We use the velocity arc ($\rho_{\text{velocity}}$) as the single-signal
classifier here because it is available for all 239 qualifying
contributors; the dual-signal (complexity $\times$ velocity) definition
of Table~\ref{tab:trajectories} is retained for the main results.

For clarity we define the \textit{velocity-rising} group as
$\rho_{\text{velocity}} > +0.2$
(resolution time increasing over career: consistent with taking on
progressively harder or broader work) and the \textit{velocity-falling}
group as $\rho_{\text{velocity}} < -0.2$.
These correspond to 14.9\,\% and 12.3\,\% of Spark contributors,
22.9\,\% and 9.6\,\% of Hadoop contributors, and 11.9\,\% and 19.0\,\%
of Camel contributors respectively.

Table~\ref{tab:growth_factors} reports the median values and
$p$-values for the factors that show significant or near-significant
differences in at least one project.

\begin{table}[h]
\centering
\caption{Selected growth-predictor comparisons:
velocity-rising vs.\ other contributors, by project.
Mann-Whitney $U$ test, two-sided.}
\label{tab:growth_factors}
\begin{tabular}{lccccc}
\toprule
 & \multicolumn{2}{c}{\textbf{SPARK}} & \multicolumn{2}{c}{\textbf{HADOOP}} \\
\cmidrule(lr){2-3} \cmidrule(lr){4-5}
Factor & Rising & Other & Rising & Other & $p$ (Sp.) \\
\midrule
Days $\Delta$ (late$-$early) & $+54$ & $+8$ & $+19$ & $+6$ & $<0.001^{***}$ / 0.45 \\
Sub-task ratio & 7.7\,\% & 13.9\,\% & 1.9\,\% & 2.5\,\% & $0.019^{*}$ / 0.67 \\
First year & 2014 & 2015 & 2009 & 2010 & $0.025^{*}$ / 0.97 \\
Late career days & 71 & 41 & 43 & 58 & $0.008^{**}$ / 0.92 \\
Bug ratio & 26.5\,\% & 24.7\,\% & 18.9\,\% & 24.4\,\% & 0.87 / 0.055 \\
Career span (years) & 5 & 5 & 6 & 6 & 0.80 / 0.74 \\
$\Delta$ git coverage & 0.00 & $-0.05$ & $-0.08$ & $-0.04$ & 0.90 / 0.38 \\
\bottomrule
\end{tabular}
\\[4pt]
\footnotesize $p$-values shown as SPARK / HADOOP. $^{*}p<.05$,
$^{**}p<.01$, $^{***}p<.001$.
\end{table}

Three findings stand out.

\textbf{Spark: a clear and replicable growth signal.}
In Spark, velocity-rising contributors differ significantly from others
on three independent predictors: a large positive days delta
($p < 0.001$), lower sub-task engagement ($p = 0.019$), and an
earlier project entry date ($p = 0.025$).
The days delta result means that rising-time contributors' resolution
times increase much more steeply over their career --- consistent with
deliberate escalation into harder problems.
The sub-task result is particularly actionable: growing Spark contributors
do roughly half as many sub-tasks as their peers, suggesting that
assigning fragmented delegated work is associated with \textit{slower}
trajectory development.

\textbf{Hadoop: no statistically distinguishable signal.}
In Hadoop, not a single predictor ($n = 83$ engineers) achieves $p < 0.05$.
The velocity-rising label in Hadoop is statistically indistinguishable
from a random split.
We interpret this as a \textit{complexity-floor masking effect}:
Hadoop tickets are inherently 6--18$\times$ harder than
Camel or Spark tickets (median complexity score 0.18 vs.\ 0.03 and 0.01),
producing high within-engineer variance in resolution times and thus
noisy $\rho$ estimates that are unreliable as trajectory signals.
The Hadoop anomaly (23\,\% velocity-rising) should therefore not be
interpreted as evidence of more growth; it reflects the broader
distribution of $\rho$ values when ticket difficulty is uniformly high.

\textbf{What makes Hadoop's velocity-rising cohort distinctive.}
Comparing \textit{only} the velocity-rising engineers across projects
reveals two significant cross-project differences
(Table~\ref{tab:growth_cross}): Hadoop's velocity-rising contributors
started contributing five years earlier (median 2009 vs.\ 2014 for Spark,
$p = 0.004$) and faced substantially harder early work
(early-career median resolution time 43.6~days vs.\ 21.5~days,
$p = 0.010$).
They also touched more files per commit from the outset
(3.50 vs.\ 2.50, $p = 0.028$) and are the only group to exhibit a
\textit{positive} file-complexity arc
($\rho_{\text{files}} = +0.04$ vs.\ $-0.11$ for Spark, $p = 0.055$).

\begin{table}[h]
\centering
\caption{Cross-project comparison: velocity-rising engineers only.
Mann-Whitney $U$ (Hadoop vs.\ Spark, two-sided).}
\label{tab:growth_cross}
\begin{tabular}{lcccc}
\toprule
Factor & CAMEL & SPARK & HADOOP & $p$ (H vs.\ S) \\
\midrule
First year (median)           & 2011  & 2014   & 2009  & $0.004^{**}$ \\
Early career days (median)    & 19.5  & 21.5   & 43.6  & $0.010^{**}$ \\
Files/commit early (median)   & 2.55  & 2.50   & 3.50  & $0.028^{*}$  \\
File-complexity arc $\rho$    & $-0.16$ & $-0.11$ & $+0.04$ & 0.055 \\
\bottomrule
\end{tabular}
\end{table}

These results suggest that Hadoop's velocity-rising cohort consists
primarily of the project's founding engineers --- contributors who
designed the distributed storage architecture from 2006--2010 and
whose rising resolution times reflect genuine architectural escalation,
not trajectory noise.
They are the exception that proves the rule: when a high-complexity
floor is paired with genuine architectural ownership from day one,
a real growth signal is present.
However, this condition is not reproducible by management intervention
--- it is a historical artefact of project founding.
The replicable mechanism, visible in Spark (and partially in Camel),
is sub-task reduction and scope expansion.

\subsection{Robustness checks}
\label{sec:robustness}

\textbf{Threshold sensitivity.}
At $|\rho| = 0.15$: Growing 14.8\,\%, Coasting 24.1\,\%.
At $|\rho| = 0.25$: Growing 8.5\,\%, Coasting 16.4\,\%.
The qualitative finding (coasting group is larger than growing group
at all thresholds) is stable.

\textbf{Minimum ticket count.}
Restricting to contributors with $\geq25$ tickets ($n=201$):
Growing 12.4\,\%, Coasting 19.4\,\%.
The distribution is stable.

\textbf{Project variation and signal reliability.}
Apache Hadoop shows lower commit-message coverage (49\,\%) and a
higher velocity-rising fraction (22.9\,\%) than Spark or Camel,
but the multi-factor analysis (Section~\ref{sec:multifactor}) reveals
that this fraction carries no statistically distinguishable signal:
no career attribute separates velocity-rising from other Hadoop
contributors.
The Hadoop distribution reflects the masking effect of a uniformly
hard ticket floor rather than a genuine excess of career growth.
Conclusions about the 70-20-10 distribution are therefore drawn
primarily from the Spark and Camel data, which produce reliable signals.

% ── 6. Discussion ───────────────────────────────────────────────────
\section{Discussion}

\subsection{Implications for the 70-20-10 model}

Our results challenge the descriptive premise of the 70-20-10 model
for software engineers.
If 70\,\% of development happened automatically through on-the-job
experience, we would expect the majority of contributors to show
a positive complexity arc over their careers.
Instead, the majority (70\,\%) show no directional signal, and the
coasting group (20\,\%) actually shows a \textit{negative} trajectory
--- declining complexity combined with slowing velocity.

This does not necessarily invalidate the 70-20-10 model as a
prescriptive framework; it may be correct that 70\,\% of
\textit{effective} professional development should be experience-based.
The empirical finding is that exposure alone does not produce this
outcome: only 11\,\% of contributors showed evidence of the
experience-driven growth the model assumes.

\subsection{The coasting paradox and its management implications}

The coasting paradox --- that the most communicatively active and
visible contributors are on a declining trajectory --- has direct
implications for talent management.

If managers use communication volume, responsiveness, or
cross-team connectivity as proxies for competence or growth,
the paradox suggests they will systematically overestimate the
development trajectory of coasting contributors and underestimate
that of growing ones.

One interpretation is that coasting contributors are \textit{efficient}
at familiar problems: they can resolve simpler tickets quickly,
comment prolifically because they recognise patterns, and are
well-connected because they have been in the community long enough
to establish relationships. None of this is dysfunctional ---
but it is different from growth.

Growing contributors, by contrast, are likely spending more
cognitive effort per ticket (unfamiliar modules, harder problems)
and less time commenting. Their lower degree centrality may
reflect genuine isolation in new problem spaces, not disengagement.

\subsection{Detecting growth from data}

The velocity arc (improving speed on harder work) is detectable
6--18 months before a measurable complexity ceiling rise in our
data. This lag represents a potential early-warning window:
managers with access to Jira and git data could identify
contributors approaching a readiness threshold \textit{before}
promotion decisions, rather than inferring retrospectively from
outcomes.

The data requirements are minimal: resolution dates, assignee
fields, and changed file lists are sufficient. No custom
instrumentation, surveys, or proprietary tooling is required.

\subsection{The complexity floor and signal reliability}

A natural hypothesis is that placing engineers in a harder project
--- one with a higher baseline complexity --- would accelerate growth.
Our cross-project comparison does not support this.

Apache Hadoop has a complexity floor roughly 6--18$\times$ higher
than Camel or Spark: median ticket score 0.18 vs.\ 0.03 and 0.01.
Yet the trajectory signal for individual Hadoop engineers is
\textit{less} reliable, not more.
The uniformly hard ticket distribution compresses the variance available
to detect monotonic trends within any individual career arc.
Equivalently: when every ticket is hard, a rising trend in difficulty
is harder to observe than when there is a wide range to move across.

The implication for talent management is that a hard assignment
environment is a necessary but insufficient condition for growth.
Complexity exposure alone does not produce the signal.
What produces the signal --- and what Spark's growing cohort
demonstrates --- is \textit{selective engagement}: fewer sub-tasks
(which fragment effort and reduce ownership), and deliberate scope
expansion that is visible in the files-per-commit trajectory.

\subsection{Practical intervention: deliberate work assignment}

The three-phase escalation arc observed in growing contributors,
combined with the multi-factor predictor analysis, suggests a
tractable intervention: deliberate escalation of ticket complexity
\textit{and} sub-task management, matched to the engineer's current phase.

\begin{enumerate}
  \item \textbf{Phase 1 (repetitions):} Assign tickets within a
    narrow domain to allow pattern recognition to accumulate.
    Monitor for velocity improvement.
    \textit{Key signal: sub-task ratio should be below 10\,\%};
    high sub-task assignment at this phase may prevent ownership
    from developing.
  \item \textbf{Phase 2 (scope expansion):} When velocity improves,
    gradually introduce cross-module tickets.
    Monitor for simultaneous velocity maintenance and
    rising files-per-commit (scope delta $> 0$).
  \item \textbf{Phase 3 (ceiling-raising):} Introduce core-module
    tickets with high watcher counts.
    These are the hardest, highest-visibility problems.
    Monitor that resolution time does not plateau --- a plateau
    with stable file scope suggests the engineer has reached the
    current ceiling and needs architectural mentorship, not
    more of the same tickets.
\end{enumerate}

All three signals --- sub-task ratio, scope delta, files-per-commit arc
--- are derivable from a team's existing Jira and git history,
requiring no additional process overhead or instrumentation.

% ── 7. Limitations ──────────────────────────────────────────────────
\section{Limitations}

\textbf{Open-source context.}
Apache contributors are volunteers with heterogeneous engagement levels.
The coasting and growing distributions may differ substantially in
commercial settings where ticket assignment is more tightly managed
and career incentives differ.
We recommend caution in generalising quantitative thresholds
(e.g.\ 11\,\% / 20\,\%) to commercial teams.

\textbf{Anonymised identities.}
Apache Jira anonymises contributor keys.
We cannot confirm that the same UUID represents the same person
across the full dataset window.
UUID stability is a property of the Jira system, but account
merges, renames, or bot accounts could introduce noise.

\textbf{Complexity as a proxy.}
Our complexity score captures breadth and core-centrality of code
changes, not semantic difficulty.
A one-line fix to a core file scores differently from a multi-file
peripheral refactor.
Secondary validation (resolution time, watcher count) supports the
score's construct validity, but a gold-standard measure of
``ticket difficulty'' does not exist in this dataset.

\textbf{Small Growing cohort.}
With 42 Growing contributors (11\,\% of 378), subgroup analyses
--- particularly the hub-file exposure effect (Section~\ref{sec:cochange})
--- are directional rather than definitive.

\textbf{Temporal confound.}
Project maturity may affect complexity distributions: as easy bugs
are resolved, the remaining open tickets in a mature project are
harder. This could inflate the apparent complexity scores for
contributors who join later.
We partially control for this by using within-contributor
rank correlation (not absolute score levels), but the confound
is not fully eliminated.

\textbf{Communication network (Spark only).}
The co-comment communication analysis was conducted on Apache Spark
only. The trajectory $\times$ communication interaction
(coasting paradox) should be replicated on Camel and Hadoop,
and ideally on commercial Jira data, before strong conclusions
are drawn.

% ── 8. Conclusion ───────────────────────────────────────────────────
\section{Conclusion}

We introduced a career arc measure derived from Jira and git
metadata and applied it to 378 contributors across three Apache
projects over 15 years.
The empirical distribution of career arcs (11\,\% Growing /
20\,\% Coasting / 70\,\% No change) directly contradicts the
descriptive assumption underlying the 70-20-10 learning framework.

A communication network analysis of Apache Spark reveals that
coasting contributors are 13$\times$ more communicatively active
and 3$\times$ more central in the co-comment graph than growing
contributors --- a pattern we call the \textit{coasting paradox}.

A multi-factor predictor analysis across 21 career variables
identifies the measurable drivers of growth.
In Spark, growing contributors are distinguished by three
independently significant factors: a steep positive days-delta
(resolution time rising as they escalate into harder work,
$p < 0.001$), below-average sub-task engagement ($p = 0.019$),
and earlier project entry ($p = 0.025$).
In Hadoop, whose ticket complexity floor is 6--18$\times$ higher
than the other two projects, no factor achieves significance:
the uniformly hard environment masks the individual career signal.
This \textit{complexity floor effect} shows that environmental
difficulty alone is insufficient to produce growth; it can, paradoxically,
make growth harder to detect and perhaps harder to achieve.

All three findings converge on the same conclusion:
experience is necessary but not sufficient for growth.
The distinguishing factors are the \textit{content} and
\textit{structure} of the work --- a low sub-task ratio,
rising scope, and deliberate escalation into harder problem spaces.
Work assignment is a learning instrument.
The signal that reveals whether it is being used effectively
is already latent in every engineering team's Jira and git history.

% ── References ──────────────────────────────────────────────────────
\section*{References}
\begin{thebibliography}{9}

\bibitem{McCall1988}
McCall, M.W., Lombardo, M.M., \& Morrison, A.M. (1988).
\textit{The Lessons of Experience: How Successful Executives Develop
on the Job}.
Lexington Books.

\bibitem{Clardy2018}
Clardy, A. (2018).
70-20-10 and the dominance of informal learning: A fact in search
of evidence.
\textit{Human Resource Development Review}, 17(2), 153--178.

\bibitem{Clauset2004}
Clauset, A., Newman, M.E.J., \& Moore, C. (2004).
Finding community structure in very large networks.
\textit{Physical Review E}, 70(6), 066111.

\bibitem{Newman2004}
Newman, M.E.J. (2004).
Analysis of weighted networks.
\textit{Physical Review E}, 70(5), 056131.

\bibitem{Fruchterman1991}
Fruchterman, T.M.J., \& Reingold, E.M. (1991).
Graph drawing by force-directed placement.
\textit{Software: Practice and Experience}, 21(11), 1129--1164.

\end{thebibliography}

\end{document}
